% answers_short.tex —— 2020 年·江苏卷 · 简短答案速查
% 由 tools/make_short_answers.py 自 answers.tex 的 \gksolhead/\gkanswer 提取，
% 并与 2020/2020_数学真题_江苏_答案要点.md 逐题比对后生成，勿手改（改源 MD 后重跑）。

\begin{gknotice}
\textbf{使用说明}：本节只给最终答案便于快速核对，完整推导见解析版。答案为编写者独立推导，\textbf{非官方标准答案}。
\end{gknotice}

\gkpart{填空题}
\begin{itemize}[leftmargin=2.6em]
  \item[\textbf{1.}] $\{0,2\}$。$A$ 与 $B$ 的公共元素只有 $0$ 与 $2$。
  \item[\textbf{2.}] $3$。$z=(1+\mathrm i)(2-\mathrm i)=2-\mathrm i+2\mathrm i-\mathrm i^2=3+\mathrm i$，实部为 $3$。
  \item[\textbf{3.}] $2$。由 $\dfrac{4+2a+3-a+5+6}5=4$ 得 $18+a=20$，故 $a=2$。
  \item[\textbf{4.}] $\dfrac19$。点数和为 $5$ 的有 $(1,4),(2,3),(3,2),(4,1)$ 共 $4$ 种，概率为 $\dfrac4{36}=\dfrac19$。
  \item[\textbf{5.}] $-3$。$x>0$ 时 $y=2^x>0$，不可能等于 $-2$；故 $x\leqslant0$，由 $x+1=-2$ 得 $x=-3$。
  \item[\textbf{6.}] $\dfrac32$。由 $\dfrac{b}{a}=\dfrac{\sqrt5}{2}$ 且 $b^2=5$ 得 $a=2$、$c=\sqrt{a^2+b^2}=3$，故 $e=\dfrac ca=\dfrac32$。
  \item[\textbf{7.}] $-4$。$f(-8)=-f(8)=-8^{\frac23}=-4$。
  \item[\textbf{8.}] $\dfrac13$。由 $\sin^2\left(\dfrac\pi4+\alpha\right)=\dfrac{1+\sin2\alpha}2=\dfrac23$ 得 $\sin2\alpha=\dfrac13$。
  \item[\textbf{9.}] $12\sqrt3-\dfrac\pi2$。底面正六边形面积 $6\times\dfrac{\sqrt3}4\times2^2=6\sqrt3$，正六棱柱体积 $6\sqrt3\times2=12\sqrt3$，内孔圆柱体积 $\pi\times0.5^2\times2=\dfrac\pi2$，故体积为 $12\sqrt3-\dfrac\pi2$（$\mathrm{cm}^3$）。
  \item[\textbf{10.}] $x=-\dfrac{5\pi}{24}$。平移后 $y=3\sin\left(2x-\dfrac\pi{12}\right)$，对称轴 $x=\dfrac{7\pi}{24}+\dfrac{k\pi}2$，$k=0$ 与 $k=-1$ 中离 $y$ 轴最近的是 $-\dfrac{5\pi}{24}$。
  \item[\textbf{11.}] $4$。由 $a_n+b_n=S_n-S_{n-1}=2n-2+2^{n-1}$（$n\geqslant1$，$S_0=0$）得 $d=2$、$q=2$，故 $d+q=4$。
  \item[\textbf{12.}] $\dfrac45$。令 $v=y^2\in(0,1]$，则 $x^2=\dfrac{1-v^2}{5v}$，$x^2+y^2=\dfrac1{5v}+\dfrac{4v}5\geqslant2\sqrt{\dfrac4{25}}=\dfrac45$，$v=\dfrac12$ 时取到。
  \item[\textbf{13.}] $\dfrac{18}5$。由 $\overrightarrow{PA}=m\overrightarrow{PB}+\left(\dfrac32-m\right)\overrightarrow{PC}$ 与 $|\overrightarrow{PA}|=9$ 解得 $m=\dfrac{27}{25}$，此时 $AD=3$、$BD=\dfrac75$，故 $CD=5-\dfrac75=\dfrac{18}5$。
  \item[\textbf{14.}] $10\sqrt5$。$PA=PB$ 使 $AB\perp PC$ 且被 $PC$ 垂直平分，设圆心到 $AB$ 的距离为 $d$，则面积 $S=\sqrt{36-d^2}\,(d+1)$，$d=4$ 时取最大值 $10\sqrt5$。
\end{itemize}

\gkpart{解答题}
\begin{itemize}[leftmargin=2.6em]
  \item[\textbf{15.}] (1) 证明。$E$，$F$ 分别是 $AC$，$B_1C$ 的中点，故 $EF\parallel AB_1$；又 $AB_1\subset$ 平面 $AB_1C_1$ 而 $EF\not\subset$ 平面 $AB_1C_1$，所以 $EF\parallel$ 平面 $AB_1C_1$；(2) 证明。$B_1C\perp$ 平面 $ABC$ 给出 $B_1C\perp AB$，又 $AB\perp AC$，故 $AB\perp$ 平面 $AB_1C$；而 $AB\subset$ 平面 $ABB_1$，所以平面 $AB_1C\perp$ 平面 $ABB_1$
  \item[\textbf{16.}] (1) $\dfrac{\sqrt5}5$。由余弦定理 $b^2=a^2+c^2-2ac\cos B=5$ 得 $b=\sqrt5$，再由正弦定理 $b\sin C=c\sin B$ 得 $\sin C=\dfrac{\sqrt5}5$；(2) $\dfrac2{11}$。$\sin\angle ADC=\dfrac35$，$\tan\angle ADC=-\dfrac34$，$\tan C=\dfrac12$，故 $\tan\angle DAC=\tan\left(180^\circ-\angle ADC-C\right)=\dfrac2{11}$
  \item[\textbf{17.}] (1) $120$ 米。$B$ 处 $b=40$ 给出 $h_2=-\dfrac{40^3}{800}+240=160$；$A$ 处 $h_1=\dfrac{a^2}{40}=160$ 得 $a=80$，故 $AB=80+40=120$（米）；(2) $O'E=20$ 米时总造价最低，最低为 $155k$ 万元。设 $O'E=x$，则 $O'C=80-x$，$CD=160-\dfrac{(80-x)^2}{40}$，$EF=160-\left(-\dfrac{x^3}{800}+6x\right)$，总造价 $W(x)=k\left[\dfrac{x^3}{800}-\dfrac{3x^2}{80}+160\right]$，$W'(x)=\dfrac{3x(x-20)}{800}k$，故 $x=20$ 时取最小值
  \item[\textbf{18.}] (1) $6$。$A\left(1,\dfrac32\right)$，周长 $=AF_1+AF_2+F_1F_2=2a+2c=4+2=6$；(2) $-4$。设 $P(p,0)$，则 $Q\left(4,\dfrac{3(p-4)}{2(p-1)}\right)$，$\overrightarrow{OP}\cdot\overrightarrow{QP}=p(p-4)=(p-2)^2-4\geqslant-4$；(3) $(2,0)$ 或 $\left(-\dfrac27,-\dfrac{12}7\right)$。$S_2=3S_1$ 等价于 $M$ 到直线 $AB:3x-4y+3=0$ 的距离为 $\dfrac95$，即 $3x-4y=6$ 或 $3x-4y=-12$，与椭圆方程联立得前一种情形两组解、后一种无解
  \item[\textbf{19.}] (1) $h(x)=2x$。由 $x^2+2x\geqslant kx+b\geqslant-x^2+2x$ 对一切 $x$ 成立，两个判别式迫使 $(k-2)^2=0$、$b=0$；(2) $[0,3]$。$f(x)\geqslant h(x)$ 等价于 $x^2-(k+1)x+(k+1)\geqslant0$，得 $-1\leqslant k\leqslant3$；$h(x)\geqslant g(x)$ 等价于 $k(x-1-\ln x)\geqslant0$，得 $k\geqslant0$；合并为 $0\leqslant k\leqslant3$；(3) 证明。$f(x)-h(x)=(x-t)^2\left(x^2+2tx+3t^2-2\right)$，$h(x)-g(x)=L^2-4(x-c)^2$（$c=\dfrac{t^3-t}2$，$L^2=t^6-5t^4+3t^2+8$）。$|t|\geqslant1$ 时 $n-m\leqslant L\leqslant\sqrt7$；$0<|t|<1$ 时 $n-m\leqslant\sqrt2+|t|<\sqrt2+1<\sqrt7$
  \item[\textbf{20.}] (1) $\lambda=1$。$k=1$ 时条件即 $a_{n+1}=\lambda a_{n+1}$ 对一切 $n$ 成立，等差数列中不可能除首项外全为 $0$，故 $\lambda=1$；(2) $a_n=\begin{cases}1,&n=1,\\3\cdot4^{n-2},&n\geqslant2.\end{cases}$ 令 $x_n=\sqrt{S_n}$，条件化为 $x_{n+1}-x_n=\dfrac{\sqrt3}3\sqrt{a_{n+1}}$，平方并消去 $a_{n+1}=x_{n+1}^2-x_n^2$ 得 $x_{n+1}=2x_n$，故 $S_n=4^{n-1}$；(3) 存在，$\lambda$ 的取值范围是 $(0,1)$。令 $x_n=S_n^{\frac13}$，条件化为 $(x_{n+1}-x_n)^3=\lambda^3\left(x_{n+1}^3-x_n^3\right)$；当 $a_{n+1}>0$ 时比值 $r=\dfrac{x_{n+1}}{x_n}$ 满足 $(r-1)^2=\lambda^3(r^2+r+1)$，该方程有解 $r>1$ 当且仅当 $0<\lambda<1$，此时按「哪些步增长」可构造出无穷多个不同数列
\end{itemize}

\gkpart{选做题}
\begin{itemize}[leftmargin=2.6em]
  \item[\textbf{21.}] 【A】（1）$a=2$，$b=2$；（2）$\boldsymbol M^{-1}=\begin{bmatrix}\dfrac25&-\dfrac15\\\dfrac15&\dfrac25\end{bmatrix}$。【B】（1）$\rho_1=4$，$\rho_2=2$；（2）公共点的极坐标为 $\left(2\sqrt2,\dfrac\pi4\right)$。【C】$\left[-2,\dfrac23\right]$。
\end{itemize}

\gkpart{解答题2}
\begin{itemize}[leftmargin=2.6em]
  \item[\textbf{22.}] (1) $\dfrac{\sqrt{15}}{15}$。以 $O$ 为原点、$BD$ 为 $x$ 轴建系，$\overrightarrow{AB}=(-1,0,-2)$，$\overrightarrow{DE}=(-1,1,1)$，$|\cos\langle\overrightarrow{AB},\overrightarrow{DE}\rangle|=\dfrac1{\sqrt{15}}=\dfrac{\sqrt{15}}{15}$；(2) $\dfrac{2\sqrt{39}}{13}$。平面 $CDE$ 的法向量取 $(2,1,1)$，平面 $FDE$ 的法向量取 $(2,7,-5)$，两平面夹角余弦为 $\dfrac1{\sqrt{13}}$，故 $\sin\theta=\sqrt{1-\dfrac1{13}}=\dfrac{2\sqrt{39}}{13}$
  \item[\textbf{23.}] (1) $p_1=\dfrac13$，$q_1=\dfrac23$，$p_2=\dfrac7{27}$，$q_2=\dfrac{16}{27}$。第一次操作中甲给出黑球（概率 $\dfrac23$）则黑球数变为 $1$，给出白球（概率 $\dfrac13$）则仍为 $2$；第二次按状态转移逐项相加；(2) $2p_n+q_n=\dfrac13\left(2p_{n-1}+q_{n-1}\right)+\dfrac23$，$E(X_n)=1+\dfrac1{3^n}$。由状态转移得 $p_n=\dfrac13p_{n-1}+\dfrac29q_{n-1}$ 与 $q_n=\dfrac23p_{n-1}+\dfrac59q_{n-1}+\dfrac23r_{n-1}$（$r_{n-1}=1-p_{n-1}-q_{n-1}$），相加即得递推式；又 $E(X_n)=2p_n+q_n$，解递推得 $E(X_n)=1+\dfrac1{3^n}$
\end{itemize}
